% JumpHMM slide (spec, slide 4): the generator's architecture across the slide (user, 2026-10-29: no caption text),
% and the stylized facts against real SPY as three panels side by side, after Fig. 3 of the JumpHMM paper. The data
% (talk/data/sf_*.csv): the generator's density sits on SPY's; its tail quantiles are as far out as SPY's or further;
% its clustering (ACF of |g|) falls from 0.22 at lag 1 to 0.03 by lag 3, while SPY's is still 0.13 at lag 26.
% the paper reference sits at the lower left, on the page number's line, after the generator's role (review round,
% 2026-10-30: the title's tail-aware forecasting is the test market, while the controller forecasts with GBM)
{\renewcommand{\vnfootleft}{The controller forecasts with GBM; JumpHMM generates the test markets. Alswaidan and
  Varner, arXiv:2603.10202.}
\begin{frame}{JumpHMM: a market generator with fat tails}
  \vnjumphmmarchitecture\par
  \vspace{0pt}
  % legend for the three panels, in their own line, band and dash styles
  {\centering\fontsize{7.5}{9}\selectfont\color{vnslidebody}%
    \tikz[baseline=-0.6ex]\draw[draw=vnfigaccent, line width=1.2pt] (0,0) -- (0.45,0);~real SPY, 2014–2024\hspace{1.5em}%
    \tikz[baseline=-0.6ex]{\fill[vnband, fill opacity=0.3] (0,-0.08) rectangle (0.45,0.08);
      \draw[draw=vnslidelink, line width=0.9pt] (0,0) -- (0.45,0);}~generator, 100 paths: median, 10–90\%\hspace{1.5em}%
    \tikz[baseline=-0.6ex]\draw[draw=vndottick, line width=0.8pt, dash pattern=on 2.5pt off 1.5pt] (0,0) --
      (0.45,0);~normal, SPY's mean and SD\par}
  % half a point more than before: the footer's descenders deepen the page number's line
  \vspace{-2.5pt}
  \vnstylizedfacts{3.35cm}{1.47cm}
  % speaker notes (user, 2026-10-31: algorithm cheat sheets), two pages. Sources: arxiv-paper/sections/
  % method_synthetic.tex (verified against code/src/SyntheticCalibration.jl SYNTH_HMM and the pinned JumpHMM.jl,
  % src/JumpHMM.jl fit and src/Tune.jl w_kappa = 0.20, acf_lags = 25); results/synthetic/generator.jld2;
  % results/synthetic/paths.csv (has_burst 523 of 10,000); results/synthetic/checks.csv (V4: -0.097, -0.171, -0.097).
  \note{\footnotesize
    \vnnotehead{JumpHMM: one series}
    \vnnoterow{States}{50 intervals of the daily growth rate, bounded by equal-probability quantiles of a Laplace
      distribution fitted to the series.}
    \vnnoterow{Transitions}{Estimated by counting the state-to-state moves in the training data.}
    \vnnoterow{Daily draw}{Step the chain, then draw a Student-$t$ growth rate ($\nu = 5$) located and scaled by the
      training days in the new state.}
    \vnnoterow{Why jumps}{The counted chain mixes within a few days, so on its own the volatility clustering would
      fade too fast.}
    \vnnoterow{Jump episodes}{On a day with no episode under way, one starts with probability ε and lasts a Poisson
      number of days with mean λ. Each episode day draws its state uniformly from the 5 most negative states with
      probability 0.52, otherwise from the 5 most positive: a sustained burst of large moves of both signs, not a
      single price jump.}
    \vnnoterow{Tuning ε and λ}{Grid search over simulated paths that contain an episode, minimizing the squared error
      of the absolute-growth autocorrelation at lags 1 to 25 plus 0.20 times the squared error of excess kurtosis.}}
\end{frame}
\note{\footnotesize
  \vnnotehead{JumpHMM in this study}
  \vnnoterow{Market model}{Fitted to SPY's closes over 2014–2024: $\epsilon = 10^{-4}$ per day, $\lambda = 118.9$
    trading days. 523 of 10,000 simulated 420-day windows (5.2\%) contain a burst.}
  \vnnoterow{Tuned values}{ε sat at the grid's lower edge, $10^{-4}$, for all 11 sector models and 370 of the 373
    name models; λ spread over its whole grid, 10 to 120 trading days.}
  \vnnoterow{Versus the paper}{50 states instead of 100, so each tail set holds 10\% of the states rather than 5\%;
    every sector and name model is tuned on its own series instead of reusing SPY's ε and λ.}
  \vnnoterow{Known limit}{Composed names cluster less than real names: the absolute-growth autocorrelation is below
    the real one by a median of 0.10, 0.17 and 0.10 at lags 1, 5 and 20.}
  \vnnoterow{Software}{JumpHMM.jl 0.1.0 at a pinned commit, unmodified; only its fit, tune and simulate functions.
    The multi-asset composition is written in the study code (next slide).}
  \vnnoterow{Role}{JumpHMM plays the market. The controller still forecasts with GBM and never sees the
    generator.}}
}
